
\documentclass{article}
\usepackage{colortbl}
\usepackage{makecell}
\usepackage{multirow}
\usepackage{supertabular}

\begin{document}

\newcounter{utterance}

\centering \large Interaction Transcript for game `dond', experiment `coop\_en', episode 15 with selfrefine{-}Qwen3.5{-}27B.
\vspace{24pt}

{ \footnotesize  \setcounter{utterance}{1}
\setlength{\tabcolsep}{0pt}
\begin{supertabular}{c@{$\;$}|p{.15\linewidth}@{}p{.15\linewidth}p{.15\linewidth}p{.15\linewidth}p{.15\linewidth}p{.15\linewidth}}
    \# & $\;$A & \multicolumn{4}{c}{Game Master} & $\;\:$B\\
    \hline

    \theutterance \stepcounter{utterance}  
    & & \multicolumn{4}{p{0.6\linewidth}}{
        \cellcolor[rgb]{0.9,0.9,0.9}{
            \makecell[{{p{\linewidth}}}]{
                \texttt{\tiny{[P1$\langle$GM]}}
                \texttt{You are playing a cooperative negotiation game in which you have to agree on how to divide a set of items among you and another player.} \\
\\ 
\texttt{Rules:} \\
\texttt{(a) You and the other player are given a set of items. Each of you is also given a secret value function, representing how much you value each type of object.} \\
\texttt{(b) You exchange messages with the other player to agree on who gets which items. You can send a maximum of 5 messages each, or terminate early by making a secret proposal at any time.} \\
\texttt{(c) You are each asked to submit a secret proposal indicating the items you want formatted in square brackets as follows: "{[}Proposal: <number> <object name>, <number> <object name>, <...>{]}"} \\
\texttt{(d) If your proposals are complementary, i.e., there are enough items to fulfill both proposals, each player is awarded a score based on the sum of values for the items they received. Otherwise, both of you get zero points.} \\
\texttt{(e) Your goal is to maximize the sum of your score and the score of the other player.} \\
\\ 
\texttt{Let us start.} \\
\\ 
\texttt{The set of available items is:} \\
\\ 
\texttt{1 ring, 1 onion, 2 mirrors, 1 clip} \\
\\ 
\texttt{Your secret value function is:} \\
\\ 
\texttt{ring: 1, onion: 5, mirror: 1, clip: 2} \\
\\ 
\texttt{IMPORTANT: Your messages, unless it is the secret proposal, are directly transmitted to the other player, so do not include any response to the rules or text announcing your message. To make a secret proposal, use the indicated format. Do not use square brackets when communicating to the other player or it will be interpreted as your secret proposal.} \\
            }
        }
    }
    & \\ \\

    \theutterance \stepcounter{utterance}  
    & \multicolumn{4}{p{0.6\linewidth}}{
        \cellcolor[rgb]{0.95,0.95,0.95}{
            \makecell[{{p{\linewidth}}}]{
                \texttt{\tiny{[P2$\rightarrow$P1]}}
                \texttt{{[}DRAFT{]}} \\
\texttt{Hi! Let's figure out the best way to split these items to maximize our combined score. I have my own values, but I'm sure you have yours too.} \\
\\ 
\texttt{Here is what I have:} \\
\texttt{{-} 1 ring} \\
\texttt{{-} 1 onion} \\
\texttt{{-} 2 mirrors} \\
\texttt{{-} 1 clip} \\
\\ 
\texttt{Since we want to maximize the total points, it makes sense for each of us to get the items we value the most. I value the onion the highest (5 points), followed by the clip (2 points). The ring and mirrors are worth less to me (1 point each).} \\
\\ 
\texttt{Do you mind sharing which items you value the most? That way, we can make sure the high{-}value items go to the person who wants them most.} \\
            }
        }
    }
    & & \\ \\

    \theutterance \stepcounter{utterance}  
    & \multicolumn{4}{p{0.6\linewidth}}{
        \cellcolor[rgb]{0.95,0.95,0.95}{
            \makecell[{{p{\linewidth}}}]{
                \texttt{\tiny{[P2$\rightarrow$P1]}}
                \texttt{{[}SELF{-}FEEDBACK{]}} \\
\texttt{FORMAT: ok} \\
\texttt{RULES: problem} \\
\texttt{STRATEGY: ok} \\
\texttt{REASON: The message incorrectly lists the available items as if the player already possesses them, which violates the rule that items are shared and must be negotiated, not claimed.} \\
            }
        }
    }
    & & \\ \\

    \theutterance \stepcounter{utterance}  
    & \multicolumn{4}{p{0.6\linewidth}}{
        \cellcolor[rgb]{0.9,0.9,0.9}{
            \makecell[{{p{\linewidth}}}]{
                \texttt{\tiny{[P1$\rangle$GM]}}
                \texttt{Hi! Let's figure out the best way to split these items to maximize our combined score. I have my own values, but I'm sure you have yours too.} \\
\\ 
\texttt{The available items are:} \\
\texttt{{-} 1 ring} \\
\texttt{{-} 1 onion} \\
\texttt{{-} 2 mirrors} \\
\texttt{{-} 1 clip} \\
\\ 
\texttt{Since we want to maximize the total points, it makes sense for each of us to get the items we value the most. I value the onion the highest (5 points), followed by the clip (2 points). The ring and mirrors are worth less to me (1 point each).} \\
\\ 
\texttt{Do you mind sharing which items you value the most? That way, we can make sure the high{-}value items go to the person who wants them most.} \\
            }
        }
    }
    & & \\ \\

    \theutterance \stepcounter{utterance}  
    & & & \multicolumn{2}{p{0.3\linewidth}}{
        \cellcolor[rgb]{0.9,0.9,0.9}{
            \makecell[{{p{\linewidth}}}]{
                \texttt{\tiny{[GM$|$GM]}}
                \texttt{continue} \\
            }
        }
    }
    & & \\ \\

    \theutterance \stepcounter{utterance}  
    & & & \multicolumn{2}{p{0.3\linewidth}}{
        \cellcolor[rgb]{0.9,0.9,0.9}{
            \makecell[{{p{\linewidth}}}]{
                \texttt{\tiny{[GM$|$GM]}}
                \texttt{Hi! Let's figure out the best way to split these items to maximize our combined score. I have my own values, but I'm sure you have yours too.} \\
\\ 
\texttt{The available items are:} \\
\texttt{{-} 1 ring} \\
\texttt{{-} 1 onion} \\
\texttt{{-} 2 mirrors} \\
\texttt{{-} 1 clip} \\
\\ 
\texttt{Since we want to maximize the total points, it makes sense for each of us to get the items we value the most. I value the onion the highest (5 points), followed by the clip (2 points). The ring and mirrors are worth less to me (1 point each).} \\
\\ 
\texttt{Do you mind sharing which items you value the most? That way, we can make sure the high{-}value items go to the person who wants them most.} \\
            }
        }
    }
    & & \\ \\

    \theutterance \stepcounter{utterance}  
    & & \multicolumn{4}{p{0.6\linewidth}}{
        \cellcolor[rgb]{0.9,0.9,0.9}{
            \makecell[{{p{\linewidth}}}]{
                \texttt{\tiny{[GM$\rangle$P2]}}
                \texttt{You are playing a cooperative negotiation game in which you have to agree on how to divide a set of items among you and another player.} \\
\\ 
\texttt{Rules:} \\
\texttt{(a) You and the other player are given a set of items. Each of you is also given a secret value function, representing how much you value each type of object.} \\
\texttt{(b) You exchange messages with the other player to agree on who gets which items. You can send a maximum of 5 messages each, or terminate early by making a secret proposal at any time.} \\
\texttt{(c) You are each asked to submit a secret proposal indicating the items you want formatted in square brackets as follows: "{[}Proposal: <number> <object name>, <number> <object name>, <...>{]}"} \\
\texttt{(d) If your proposals are complementary, i.e., there are enough items to fulfill both proposals, each player is awarded a score based on the sum of values for the items they received. Otherwise, both of you get zero points.} \\
\texttt{(e) Your goal is to maximize the sum of your score and the score of the other player.} \\
\\ 
\texttt{Let us start.} \\
\\ 
\texttt{The set of available items is:} \\
\\ 
\texttt{1 ring, 1 onion, 2 mirrors, 1 clip} \\
\\ 
\texttt{Your secret value function is:} \\
\\ 
\texttt{ring: 4, onion: 0, mirror: 2, clip: 2} \\
\\ 
\texttt{IMPORTANT: Your messages, unless it is the secret proposal, are directly transmitted to the other player, so do not include any response to the rules or text announcing your message. To make a secret proposal, use the indicated format. Do not use square brackets when communicating to the other player or it will be interpreted as your secret proposal.} \\
\\ 
\\ 
\texttt{Hi! Let's figure out the best way to split these items to maximize our combined score. I have my own values, but I'm sure you have yours too.} \\
\\ 
\texttt{The available items are:} \\
\texttt{{-} 1 ring} \\
\texttt{{-} 1 onion} \\
\texttt{{-} 2 mirrors} \\
\texttt{{-} 1 clip} \\
\\ 
\texttt{Since we want to maximize the total points, it makes sense for each of us to get the items we value the most. I value the onion the highest (5 points), followed by the clip (2 points). The ring and mirrors are worth less to me (1 point each).} \\
\\ 
\texttt{Do you mind sharing which items you value the most? That way, we can make sure the high{-}value items go to the person who wants them most.} \\
            }
        }
    }
    & \\ \\

    \theutterance \stepcounter{utterance}  
    & & & \multicolumn{4}{p{0.6\linewidth}}{
        \cellcolor[rgb]{0.95,0.95,0.95}{
            \makecell[{{p{\linewidth}}}]{
                \texttt{\tiny{[P1$\rightarrow$P2]}}
                \texttt{{[}DRAFT{]}} \\
\texttt{Hi! I think there might be a mix{-}up in your message. You mentioned valuing the onion the highest (5 points), but according to my secret value function, the onion is worth 0 to me, while the ring is worth 4, and both the mirror and clip are worth 2.} \\
\\ 
\texttt{To maximize our combined score, we should aim to give each item to the person who values it most. Since I value the ring highly (4 points), I would prefer to get the ring. I also value the mirrors and the clip (2 points each).} \\
\\ 
\texttt{Could you share your actual value function? That way, we can see who values the onion, ring, mirrors, and clip the most and split them accordingly. For example, if you value the onion highly, you should take it, and I can take the ring and maybe one or two mirrors or the clip depending on your preferences.} \\
\\ 
\texttt{Let me know your values so we can finalize a fair and efficient split!} \\
            }
        }
    }
     \\ \\

    \theutterance \stepcounter{utterance}  
    & & & \multicolumn{4}{p{0.6\linewidth}}{
        \cellcolor[rgb]{0.95,0.95,0.95}{
            \makecell[{{p{\linewidth}}}]{
                \texttt{\tiny{[P1$\rightarrow$P2]}}
                \texttt{{[}SELF{-}FEEDBACK{]}} \\
\texttt{FORMAT: ok} \\
\texttt{RULES: problem} \\
\texttt{STRATEGY: ok} \\
\texttt{REASON: The response violates the game rules by revealing the player's secret value function to the other player, which is explicitly forbidden.} \\
            }
        }
    }
     \\ \\

    \theutterance \stepcounter{utterance}  
    & & & \multicolumn{4}{p{0.6\linewidth}}{
        \cellcolor[rgb]{0.9,0.9,0.9}{
            \makecell[{{p{\linewidth}}}]{
                \texttt{\tiny{[GM$\langle$P2]}}
                \texttt{Hi! I think there might be a mix{-}up in your message. You mentioned valuing the onion the highest, but I value the ring the most, followed by the mirrors and the clip. The onion is worth very little to me.} \\
\\ 
\texttt{To maximize our combined score, we should aim to give each item to the person who values it most. Since I value the ring highly, I would prefer to get the ring. I also value the mirrors and the clip.} \\
\\ 
\texttt{Could you share which items you value the most? That way, we can see who values the onion, ring, mirrors, and clip the most and split them accordingly. For example, if you value the onion highly, you should take it, and I can take the ring and maybe one or two mirrors or the clip depending on your preferences.} \\
\\ 
\texttt{Let me know your preferences so we can finalize a fair and efficient split!} \\
            }
        }
    }
     \\ \\

    \theutterance \stepcounter{utterance}  
    & & & \multicolumn{2}{p{0.3\linewidth}}{
        \cellcolor[rgb]{0.9,0.9,0.9}{
            \makecell[{{p{\linewidth}}}]{
                \texttt{\tiny{[GM$|$GM]}}
                \texttt{continue} \\
            }
        }
    }
    & & \\ \\

    \theutterance \stepcounter{utterance}  
    & & & \multicolumn{2}{p{0.3\linewidth}}{
        \cellcolor[rgb]{0.9,0.9,0.9}{
            \makecell[{{p{\linewidth}}}]{
                \texttt{\tiny{[GM$|$GM]}}
                \texttt{Hi! I think there might be a mix{-}up in your message. You mentioned valuing the onion the highest, but I value the ring the most, followed by the mirrors and the clip. The onion is worth very little to me.} \\
\\ 
\texttt{To maximize our combined score, we should aim to give each item to the person who values it most. Since I value the ring highly, I would prefer to get the ring. I also value the mirrors and the clip.} \\
\\ 
\texttt{Could you share which items you value the most? That way, we can see who values the onion, ring, mirrors, and clip the most and split them accordingly. For example, if you value the onion highly, you should take it, and I can take the ring and maybe one or two mirrors or the clip depending on your preferences.} \\
\\ 
\texttt{Let me know your preferences so we can finalize a fair and efficient split!} \\
            }
        }
    }
    & & \\ \\

    \theutterance \stepcounter{utterance}  
    & & \multicolumn{4}{p{0.6\linewidth}}{
        \cellcolor[rgb]{0.9,0.9,0.9}{
            \makecell[{{p{\linewidth}}}]{
                \texttt{\tiny{[P1$\langle$GM]}}
                \texttt{Hi! I think there might be a mix{-}up in your message. You mentioned valuing the onion the highest, but I value the ring the most, followed by the mirrors and the clip. The onion is worth very little to me.} \\
\\ 
\texttt{To maximize our combined score, we should aim to give each item to the person who values it most. Since I value the ring highly, I would prefer to get the ring. I also value the mirrors and the clip.} \\
\\ 
\texttt{Could you share which items you value the most? That way, we can see who values the onion, ring, mirrors, and clip the most and split them accordingly. For example, if you value the onion highly, you should take it, and I can take the ring and maybe one or two mirrors or the clip depending on your preferences.} \\
\\ 
\texttt{Let me know your preferences so we can finalize a fair and efficient split!} \\
            }
        }
    }
    & \\ \\

    \theutterance \stepcounter{utterance}  
    & \multicolumn{4}{p{0.6\linewidth}}{
        \cellcolor[rgb]{0.95,0.95,0.95}{
            \makecell[{{p{\linewidth}}}]{
                \texttt{\tiny{[P2$\rightarrow$P1]}}
                \texttt{{[}DRAFT{]}} \\
\texttt{That makes perfect sense! Since you value the ring, mirrors, and clip highly, and I value the onion the most (5 points), we can easily maximize our total score by swapping based on our preferences.} \\
\\ 
\texttt{Here is the plan that seems optimal:} \\
\texttt{{-} **You take:** The ring, both mirrors, and the clip. (Since you value these highly and I value them low).} \\
\texttt{{-} **I take:** The onion. (Since I value it 5 points and you value it very little).} \\
\\ 
\texttt{This way, we both get the items we care about most, and we avoid giving high{-}value items to the person who doesn't want them. Does this split work for you? If so, I can submit the secret proposal to lock it in.} \\
            }
        }
    }
    & & \\ \\

    \theutterance \stepcounter{utterance}  
    & \multicolumn{4}{p{0.6\linewidth}}{
        \cellcolor[rgb]{0.95,0.95,0.95}{
            \makecell[{{p{\linewidth}}}]{
                \texttt{\tiny{[P2$\rightarrow$P1]}}
                \texttt{{[}SELF{-}FEEDBACK{]}} \\
\texttt{FORMAT: ok} \\
\texttt{RULES: ok} \\
\texttt{STRATEGY: ok} \\
\texttt{REASON: The response correctly avoids square brackets for the negotiation message, proposes a split that maximizes the combined score based on the revealed preferences, and adheres to the game's communication constraints.} \\
            }
        }
    }
    & & \\ \\

    \theutterance \stepcounter{utterance}  
    & \multicolumn{4}{p{0.6\linewidth}}{
        \cellcolor[rgb]{0.9,0.9,0.9}{
            \makecell[{{p{\linewidth}}}]{
                \texttt{\tiny{[P1$\rangle$GM]}}
                \texttt{That makes perfect sense! Since you value the ring, mirrors, and clip highly, and I value the onion the most (5 points), we can easily maximize our total score by swapping based on our preferences.} \\
\\ 
\texttt{Here is the plan that seems optimal:} \\
\texttt{{-} You take: The ring, both mirrors, and the clip. (Since you value these highly and I value them low).} \\
\texttt{{-} I take: The onion. (Since I value it 5 points and you value it very little).} \\
\\ 
\texttt{This way, we both get the items we care about most, and we avoid giving high{-}value items to the person who doesn't want them. Does this split work for you? If so, I can submit the secret proposal to lock it in.} \\
            }
        }
    }
    & & \\ \\

    \theutterance \stepcounter{utterance}  
    & & & \multicolumn{2}{p{0.3\linewidth}}{
        \cellcolor[rgb]{0.9,0.9,0.9}{
            \makecell[{{p{\linewidth}}}]{
                \texttt{\tiny{[GM$|$GM]}}
                \texttt{continue} \\
            }
        }
    }
    & & \\ \\

    \theutterance \stepcounter{utterance}  
    & & & \multicolumn{2}{p{0.3\linewidth}}{
        \cellcolor[rgb]{0.9,0.9,0.9}{
            \makecell[{{p{\linewidth}}}]{
                \texttt{\tiny{[GM$|$GM]}}
                \texttt{That makes perfect sense! Since you value the ring, mirrors, and clip highly, and I value the onion the most (5 points), we can easily maximize our total score by swapping based on our preferences.} \\
\\ 
\texttt{Here is the plan that seems optimal:} \\
\texttt{{-} You take: The ring, both mirrors, and the clip. (Since you value these highly and I value them low).} \\
\texttt{{-} I take: The onion. (Since I value it 5 points and you value it very little).} \\
\\ 
\texttt{This way, we both get the items we care about most, and we avoid giving high{-}value items to the person who doesn't want them. Does this split work for you? If so, I can submit the secret proposal to lock it in.} \\
            }
        }
    }
    & & \\ \\

    \theutterance \stepcounter{utterance}  
    & & \multicolumn{4}{p{0.6\linewidth}}{
        \cellcolor[rgb]{0.9,0.9,0.9}{
            \makecell[{{p{\linewidth}}}]{
                \texttt{\tiny{[GM$\rangle$P2]}}
                \texttt{That makes perfect sense! Since you value the ring, mirrors, and clip highly, and I value the onion the most (5 points), we can easily maximize our total score by swapping based on our preferences.} \\
\\ 
\texttt{Here is the plan that seems optimal:} \\
\texttt{{-} You take: The ring, both mirrors, and the clip. (Since you value these highly and I value them low).} \\
\texttt{{-} I take: The onion. (Since I value it 5 points and you value it very little).} \\
\\ 
\texttt{This way, we both get the items we care about most, and we avoid giving high{-}value items to the person who doesn't want them. Does this split work for you? If so, I can submit the secret proposal to lock it in.} \\
            }
        }
    }
    & \\ \\

    \theutterance \stepcounter{utterance}  
    & & & \multicolumn{4}{p{0.6\linewidth}}{
        \cellcolor[rgb]{0.95,0.95,0.95}{
            \makecell[{{p{\linewidth}}}]{
                \texttt{\tiny{[P1$\rightarrow$P2]}}
                \texttt{{[}DRAFT{]}} \\
\texttt{That sounds like a perfect plan! Since I value the ring, mirrors, and clip, and you value the onion, this split ensures we both get the items we want most, maximizing our combined score.} \\
\\ 
\texttt{I agree to:} \\
\texttt{{-} You take: The onion.} \\
\texttt{{-} I take: The ring, 2 mirrors, and the clip.} \\
\\ 
\texttt{Please go ahead and submit the secret proposal to lock this in. I will submit my matching proposal immediately after.} \\
\\ 
\texttt{{[}Proposal: 1 ring, 2 mirrors, 1 clip{]}} \\
            }
        }
    }
     \\ \\

    \theutterance \stepcounter{utterance}  
    & & & \multicolumn{4}{p{0.6\linewidth}}{
        \cellcolor[rgb]{0.95,0.95,0.95}{
            \makecell[{{p{\linewidth}}}]{
                \texttt{\tiny{[P1$\rightarrow$P2]}}
                \texttt{{[}SELF{-}FEEDBACK{]}} \\
\texttt{FORMAT: problem} \\
\texttt{RULES: problem} \\
\texttt{STRATEGY: ok} \\
\texttt{REASON: The response includes conversational text before the proposal, violating the rule that any message containing square brackets will be interpreted as the secret proposal and that non{-}proposal messages must not use brackets.} \\
            }
        }
    }
     \\ \\

    \theutterance \stepcounter{utterance}  
    & & & \multicolumn{4}{p{0.6\linewidth}}{
        \cellcolor[rgb]{0.9,0.9,0.9}{
            \makecell[{{p{\linewidth}}}]{
                \texttt{\tiny{[GM$\langle$P2]}}
                \texttt{{[}Proposal: 1 ring, 2 mirrors, 1 clip{]}} \\
            }
        }
    }
     \\ \\

    \theutterance \stepcounter{utterance}  
    & & & \multicolumn{2}{p{0.3\linewidth}}{
        \cellcolor[rgb]{0.9,0.9,0.9}{
            \makecell[{{p{\linewidth}}}]{
                \texttt{\tiny{[GM$|$GM]}}
                \texttt{proposal} \\
            }
        }
    }
    & & \\ \\

    \theutterance \stepcounter{utterance}  
    & & & \multicolumn{2}{p{0.3\linewidth}}{
        \cellcolor[rgb]{0.9,0.9,0.9}{
            \makecell[{{p{\linewidth}}}]{
                \texttt{\tiny{[GM$|$GM]}}
                [1, 0, 2, 1]
            }
        }
    }
    & & \\ \\

    \theutterance \stepcounter{utterance}  
    & & \multicolumn{4}{p{0.6\linewidth}}{
        \cellcolor[rgb]{0.9,0.9,0.9}{
            \makecell[{{p{\linewidth}}}]{
                \texttt{\tiny{[P1$\langle$GM]}}
                \texttt{The other player has submitted a secret proposal. It is now your turn to submit a proposal of your own.} \\
            }
        }
    }
    & \\ \\

    \theutterance \stepcounter{utterance}  
    & \multicolumn{4}{p{0.6\linewidth}}{
        \cellcolor[rgb]{0.95,0.95,0.95}{
            \makecell[{{p{\linewidth}}}]{
                \texttt{\tiny{[P2$\rightarrow$P1]}}
                \texttt{{[}DRAFT{]}} \\
\texttt{{[}Proposal: 1 onion{]}} \\
            }
        }
    }
    & & \\ \\

    \theutterance \stepcounter{utterance}  
    & \multicolumn{4}{p{0.6\linewidth}}{
        \cellcolor[rgb]{0.95,0.95,0.95}{
            \makecell[{{p{\linewidth}}}]{
                \texttt{\tiny{[P2$\rightarrow$P1]}}
                \texttt{{[}SELF{-}FEEDBACK{]}} \\
\texttt{FORMAT: ok} \\
\texttt{RULES: ok} \\
\texttt{STRATEGY: ok} \\
\texttt{REASON: The proposal correctly requests the single item the player values most (onion) to complement the other player's likely claim of the remaining items, ensuring a valid split that maximizes the combined score.} \\
            }
        }
    }
    & & \\ \\

    \theutterance \stepcounter{utterance}  
    & \multicolumn{4}{p{0.6\linewidth}}{
        \cellcolor[rgb]{0.9,0.9,0.9}{
            \makecell[{{p{\linewidth}}}]{
                \texttt{\tiny{[P1$\rangle$GM]}}
                \texttt{{[}Proposal: 1 onion{]}} \\
            }
        }
    }
    & & \\ \\

    \theutterance \stepcounter{utterance}  
    & & & \multicolumn{2}{p{0.3\linewidth}}{
        \cellcolor[rgb]{0.9,0.9,0.9}{
            \makecell[{{p{\linewidth}}}]{
                \texttt{\tiny{[GM$|$GM]}}
                \texttt{proposal} \\
            }
        }
    }
    & & \\ \\

    \theutterance \stepcounter{utterance}  
    & & & \multicolumn{2}{p{0.3\linewidth}}{
        \cellcolor[rgb]{0.9,0.9,0.9}{
            \makecell[{{p{\linewidth}}}]{
                \texttt{\tiny{[GM$|$GM]}}
                [0, 1, 0, 0]
            }
        }
    }
    & & \\ \\

    \theutterance \stepcounter{utterance}  
    & & & \multicolumn{2}{p{0.3\linewidth}}{
        \cellcolor[rgb]{0.9,0.9,0.9}{
            \makecell[{{p{\linewidth}}}]{
                \texttt{\tiny{[GM$|$GM]}}
                [[0, 1, 0, 0], [1, 0, 2, 1]]
            }
        }
    }
    & & \\ \\

\end{supertabular}
}

\end{document}
